% ID: SEL-D0033-OA-E02
% ESTADO: validada. S034, cotejada con examen y criterios.
La derivada es $c'(t)=3(t-3)(t-7)$.
Es positiva entre $0$ y $3$, negativa entre $3$ y $7$
y positiva entre $7$ y $12$.

\textbf{a)} Para hallar el máximo absoluto comparamos extremos
del dominio y puntos críticos:
\[
c(0)=10,\quad c(3)=91,\quad c(7)=59,\quad c(12)=334.
\]
El máximo se alcanza en $\boxed{t=12}$ y equivale a
$\boxed{334000}$ toneladas.

\textbf{b)} El consumo decrece en $\boxed{[3,7]}$.

\textbf{c)} El esbozo recorre $(0,10)$, $(3,91)$, $(7,59)$
y $(12,334)$. La segunda derivada $6t-30$ cambia de signo en
$t=5$, donde $c(5)=75$.
\begin{center}
\setlength{\unitlength}{5mm}
\begin{picture}(17,10)(-2,-1)
\qbezier(0,0)(6.5,0)(13,0)\put(13.2,0){$t$}
\qbezier(0,0)(0,4.4)(0,8.8)\put(0,8.9){$c$}
\qbezier(0,0.25000)(0.5,1.03125)(1,1.47500)
\qbezier(1,1.47500)(1.5,1.91875)(2,2.10000)
\qbezier(2,2.10000)(2.5,2.28125)(3,2.27500)
\qbezier(3,2.27500)(3.5,2.26875)(4,2.15000)
\qbezier(4,2.15000)(4.5,2.03125)(5,1.87500)
\qbezier(5,1.87500)(5.5,1.71875)(6,1.60000)
\qbezier(6,1.60000)(6.5,1.48125)(7,1.47500)
\qbezier(7,1.47500)(7.5,1.46875)(8,1.65000)
\qbezier(8,1.65000)(8.5,1.83125)(9,2.27500)
\qbezier(9,2.27500)(9.5,2.71875)(10,3.50000)
\qbezier(10,3.50000)(10.5,4.28125)(11,5.47500)
\qbezier(11,5.47500)(11.5,6.66875)(12,8.35000)
\put(3,2.5){\makebox(0,0)[b]{$(3,91)$}}
\put(7,1.3){\makebox(0,0)[t]{$(7,59)$}}
\put(12,8.4){$(12,334)$}
\end{picture}
\end{center}
